\documentclass[11pt]{article}
\usepackage[T1]{fontenc}
\usepackage[utf8]{inputenc}
\usepackage{lmodern}
\usepackage[a4paper,margin=27mm]{geometry}
\usepackage{amsmath,amssymb,amsthm,booktabs}
\usepackage[colorlinks=true,allcolors=blue]{hyperref}
\newcommand{\Tr}{\operatorname{Tr}}
\newcommand{\id}{\operatorname{id}}
\newcommand{\I}{\mathbb{I}}
\newcommand{\ket}[1]{|#1\rangle}
\newcommand{\bra}[1]{\langle#1|}
\newtheorem{theorem}{Theorem}
\title{Exact conditional min-entropy of two-qubit unitary channels with white noise}
\author{Xiaojian Zeng\\
\small PBC School of Finance, Tsinghua University, 100190\\
\small \texttt{xiaojianzeng@foxmail.com}}
\date{}
\begin{document}
\maketitle
\begin{abstract}
The conditional min-entropy of a bipartite channel involves a
trace-preserving marginal constraint absent from the corresponding state
optimization. For an arbitrary two-qubit unitary mixed with the maximally
mixed replacer, we prove that this constraint does not change the optimum.
If $\rho_U$ is the conditioning-side reduction of the normalized pure
unitary Choi state and $L=\Tr\sqrt{\rho_U}$, both semidefinite objectives
equal $(1-p)L^2+p/4$, where $p$ is the replacer weight.
A common optimizer is $(1-p)L\sqrt{\rho_U}+(p/16)\I_4$.
Matching positive witnesses establish the expression even when $\rho_U$
is singular; the known square-root marginal property of locally maximally
mixed two-qubit states makes the optimizer channel-feasible.
This resolves the two-qubit white-noise extension posed by Badhani and Das
(arXiv:2604.01217v2), using their noiseless result and qubit marginal lemma.
The conclusion concerns single-use min-entropy for this specified noise
family.
\end{abstract}

\section{Introduction}
Conditional channel entropies quantify quantum processes in the presence
of a conditioning process. Their optimization differs from the state
problem because the comparison operator must correspond to a
trace-preserving channel \cite{Das2024,Badhani2026}.
For conditional min-entropy this distinction is a linear marginal
constraint in a semidefinite program (SDP).
Understanding when that constraint costs nothing allows a channel
quantity to be evaluated directly from a normalized Choi state.

Badhani and Das prove equality for noiseless two-qubit unitaries in
Proposition~10 of Ref.~\cite{Badhani2026}.
Their Figure~2 reports the same behavior for identity, CNOT and SWAP
with white noise, and the following paragraph asks whether it holds for
arbitrary noisy two-qubit unitaries.
We establish that two-qubit extension and give the exact objective and
a common optimizer for every noise weight.
The essential ingredients are an explicit primal--dual certificate for
a pure projector plus scalar noise, and the qubit square-root marginal
lemma already given in Appendix~D of Ref.~\cite{Badhani2026}.
The pure-state min-entropy formula \cite{KRS2009}, qubit normal forms and
SDP averaging are established tools; no new general symmetry principle is
needed.

The statement of the open extension persists in version~2 of the source.
A bounded primary-literature comparison through 12 October 2026 found
no later resolution of this exact channel identity. This comparison does
not establish exhaustive priority for the intermediate state formula.
The tele-covariant von Neumann entropy identity in
Ref.~\cite{Das2024} is a different result, and its examples mixing a gate
with the identity channel use a different noise map.

\section{Model and normalization}
All four ports $A',B',A,B$ are qubits. For $U\in U(4)$ and $0\leq p\leq1$,
define the completely positive, trace-preserving map
\begin{equation}
 \mathcal N_{p,U}(Z)=(1-p)UZU^\dagger
                  +p\,\Tr(Z)\frac{\I_{AB}}4 .
 \label{eq:channel}
\end{equation}
Thus $p$ denotes the white-noise weight, not the unitary weight.
Let $\Phi_2$ be the projector on
$(\ket{00}+\ket{11})/\sqrt2$, and set
\begin{equation}
 \Omega=(\id_{R_AR_B}\otimes\mathcal N_{p,U})
       (\Phi_{R_AA'}\otimes\Phi_{R_BB'}).
\end{equation}
We group factors as $X=R_AA$ and $Y=R_BB$, each of dimension four.
With this ordering,
\begin{equation}
 \Omega=(1-p)\ket u\bra u+\frac{p}{16}\I_{XY},\qquad
 u_{r_a,a,r_b,b}=\frac12 U_{2a+b,\,2r_a+r_b}.
 \label{eq:choi}
\end{equation}
Here $r_a,a,r_b,b\in\{0,1\}$ and $\bra u u\rangle=1$.
In particular $\Tr\Omega=1$.
Define the reduction of the \emph{pure unitary} Choi state by
\begin{equation}
 \rho_U=\Tr_X\ket u\bra u,\qquad L=\Tr\sqrt{\rho_U}.
 \label{eq:rho}
\end{equation}
The distinction between $\rho_U$ and $\Tr_X\Omega$ is essential.

The two objectives under comparison are
\begin{align}
 s(\Omega)&=\min_{Q\succeq0}
       \{\Tr Q:\I_X\otimes Q\succeq\Omega\}, \label{eq:s}\\
 t(\Omega)&=\min_{Q\succeq0}
       \left\{\Tr Q:\I_X\otimes Q\succeq\Omega,\
       \Tr_B Q=\frac{\Tr Q}{2}\I_{R_B}\right\}. \label{eq:t}
\end{align}
The state conditional min-entropy is
$H_{\min}(X|Y)_\Omega=-\log_2 s(\Omega)$
\cite{KRS2009}.
In the channel convention of Refs.~\cite{Das2024,Badhani2026},
\begin{equation}
 S_\infty[A|B]_{\mathcal N}=-\log_2 t(\Omega)-1.
 \label{eq:offset}
\end{equation}
For clarity, Eq.~(87) of Ref.~\cite{Das2024} uses the unnormalized
Choi operator $\Gamma_{\mathcal N}=4\Omega$.
Its variable $M=4Q$ has objective $\Tr M/|B'|=2\Tr Q$.
Its marginal constraint becomes exactly that in Eq.~\eqref{eq:t},
which explains the one-bit offset in Eq.~\eqref{eq:offset}.
Scalar identity operators give feasible points, and trace-bounded
positive operators form a compact set, so the minima are attained.

\section{Exact objective and common optimizer}
\begin{theorem}
For every $U\in U(4)$ and every $p\in[0,1]$, the objectives
in Eqs.~\eqref{eq:s}--\eqref{eq:t} satisfy
\begin{equation}
 s(\Omega)=t(\Omega)=(1-p)L^2+\frac p4.
 \label{eq:result}
\end{equation}
A common optimizer is
\begin{equation}
 Q_*=(1-p)L\sqrt{\rho_U}+\frac p{16}\I_Y.
 \label{eq:optimizer}
\end{equation}
Consequently,
\begin{equation}
 S_\infty[A|B]_{\mathcal N_{p,U}}
 =H_{\min}(X|Y)_\Omega-1
 =-\log_2\left((1-p)L^2+\frac p4\right)-1.
 \label{eq:entropy}
\end{equation}
The result includes zero Schmidt coefficients, repeated eigenvalues,
and both noise endpoints.
\end{theorem}

\begin{proof}
First obtain matching bounds for the state problem. Write the
positive-coefficient Schmidt decomposition as
\[
 \ket u=\sum_{i=1}^r c_i\ket{x_i}\ket{y_i},
 \qquad c_i>0,\quad \sum_i c_i^2=1.
\]
Then $L=\sum_i c_i$ and
$\sqrt{\rho_U}=\sum_i c_i\ket{y_i}\bra{y_i}$.
For an arbitrary vector $z$, weighted Cauchy--Schwarz gives
\begin{align}
 |\bra u z\rangle|^2
 &\leq L\sum_i c_i|\bra{x_i y_i}z\rangle|^2\\
 &\leq \bra z(\I_X\otimes L\sqrt{\rho_U})\ket z.
\end{align}
The second inequality includes all components along $X$ for each
$y_i$. Therefore
$\I_X\otimes L\sqrt{\rho_U}\succeq\ket u\bra u$.
Put $\alpha=1-p$ and $\beta=p/16$.
Equation~\eqref{eq:optimizer} is state-feasible because its slack is
$\alpha(\I_X\otimes L\sqrt{\rho_U}-\ket u\bra u)$,
and its trace is $\alpha L^2+4\beta$.

For a lower bound let
\begin{equation}
 \Pi=\sum_{i=1}^r\ket{y_i}\bra{y_i},\quad
 \ket w=\sum_{i=1}^r\ket{x_i}\ket{y_i},\quad
 W=\ket w\bra w+\frac{\I_X}{4}\otimes(\I_Y-\Pi).
 \label{eq:dual}
\end{equation}
We have $W\succeq0$, $\Tr_X W=\I_Y$, $\Tr W=4$, and
$\bra u W\ket u=L^2$. For every state-feasible $Q$,
\begin{equation}
 \Tr Q=\Tr[(\I_X\otimes Q)W]
       \geq\Tr(\Omega W)=\alpha L^2+4\beta.
\end{equation}
This proves $s(\Omega)=\alpha L^2+4\beta$ by weak duality.
The support-complement term in Eq.~\eqref{eq:dual} is needed for
$r<4$ and $p>0$: omitting it loses $(4-r)\beta$ from the lower bound.
No inverse or full-rank assumption has been used.

It remains to show that $Q_*$ obeys the channel marginal.
Unitarity gives both full reference and full output marginals
$\I_4/4$ for $\ket u\bra u$.
Tracing the other qubit then gives
\begin{equation}
 \Tr_B\rho_U=\frac{\I_{R_B}}2,\qquad
 \Tr_{R_B}\rho_U=\frac{\I_B}2.
 \label{eq:lmm}
\end{equation}
The square-root marginal lemma for such two-qubit states is
Appendix~D, Lemma~9 of Ref.~\cite{Badhani2026}.
For completeness, expand
\[
 \rho_U=\frac14\left(\I\otimes\I+
       \sum_{i,j=1}^3 T_{ij}\sigma_i\otimes\sigma_j\right),
 \qquad T_{ij}\in\mathbb R.
\]
There are no single-Pauli terms by Eq.~\eqref{eq:lmm}.
A real singular-value decomposition, with reflection signs absorbed
in the diagonal entries, diagonalizes $T$ using two proper rotations.
Local qubit unitaries implement these rotations.
In that local basis $\rho_U$ is a linear combination of
$\I\otimes\I$, $\sigma_x\otimes\sigma_x$,
$\sigma_y\otimes\sigma_y$ and $\sigma_z\otimes\sigma_z$.
These commute and share the four Bell vectors as an eigenbasis.
Its positive square root is therefore a sum of Bell projectors with
weights $\sqrt{\lambda_j}$. Each projector has marginal $\I_2/2$.
Returning to the original basis yields
\begin{equation}
 \Tr_B\sqrt{\rho_U}=\frac L2\I_{R_B}.
 \label{eq:sqrt}
\end{equation}
Zero or repeated eigenvalues do not affect this spectral argument.
Using $\Tr_B\I_Y=2\I_{R_B}$,
\[
 \Tr_B Q_*=\left(\frac{\alpha L^2}{2}+2\beta\right)\I_{R_B}
          =\frac{\Tr Q_*}{2}\I_{R_B}.
\]
Thus $Q_*$ is feasible for $t$. Feasible-set inclusion gives
$s\leq t\leq\Tr Q_*=s$, proving the assertion.
All inequalities remain valid for $p=0$ and $p=1$.
\end{proof}

\subsection{Why the marginal constraint has zero cost}
The same equality has a structural averaging interpretation.
Choose a local-unitary Bell eigenbasis $\{y_i\}_{i=1}^4$ of $\rho_U$
and corresponding Schmidt vectors $\{x_i\}$, completed orthonormally
when an eigenvalue is zero. For each sign vector
$\varepsilon\in\{-1,1\}^4$, define
\[
 V_X(\varepsilon)=\sum_i\varepsilon_i\ket{x_i}\bra{x_i},
 \qquad
 V_Y(\varepsilon)=\sum_i\varepsilon_i\ket{y_i}\bra{y_i}.
\]
Their tensor product fixes $\ket u$, and therefore fixes $\Omega$.
Conjugation maps every state-feasible $Q$ to another feasible operator.
Averaging over the sixteen signs preserves its trace and produces
\[
 \overline Q=\sum_i\bra{y_i}Q\ket{y_i}\ket{y_i}\bra{y_i}.
\]
Every $y_i$ has marginal $\I_2/2$, so $\overline Q$ obeys
the channel constraint. This identifies the reason that the constraint
has zero objective penalty. The averaging operations act on grouped
spaces; their use in the proof does not assert their physical
availability as local port operations.

\section{Special gates, interpretation and limits}
Let $\kappa=L^2$. Normalization and Cauchy--Schwarz give
$1\leq\kappa\leq4$. Thus the gate dependence enters only through
the Schmidt spectrum of its pure normalized Choi vector.
The quantity $2^{-H_{\min}(X|Y)_\Omega}$ is affine in the replacer
weight, although the entropy itself is logarithmic.
The identity, CNOT and SWAP have $\kappa=1,2,4$, respectively.
Their objectives and endpoints are shown in Table~\ref{tab:gates}.
These are exact instances of Eq.~\eqref{eq:result}, recovering the
three examples of Ref.~\cite{Badhani2026}.

\begin{table}[ht]
\centering
\caption{Exact common objective $s=t$ for three gates, with $p$
the maximally mixed replacer weight. Values are rational, not SDP
solver estimates. Channel entropy is $-\log_2(s)-1$.}
\label{tab:gates}
\begin{tabular}{lcccc}
\toprule
Gate & $\kappa$ & $p=0$ & $p=1/2$ & $p=1$\\
\midrule
Identity & 1 & 1 & $5/8$ & $1/4$\\
CNOT     & 2 & 2 & $9/8$ & $1/4$\\
SWAP     & 4 & 4 & $17/8$ & $1/4$\\
\bottomrule
\end{tabular}
\end{table}

For example, the noiseless SWAP has state entropy $-2$ bits and
channel entropy $-3$ bits. The completely mixed replacer has
state entropy $+2$ bits and channel entropy $+1$ bit for every $U$.
These endpoints also check the normalized Choi convention and
the one-bit channel offset.

The result removes the need to solve either SDP once the four
eigenvalues of $\rho_U$ are known. In the conditional athermality
framework of Ref.~\cite{Badhani2026}, it evaluates the single-use
channel min-entropy entering the trivial-Hamiltonian resource
description. No additional thermodynamic conversion theorem is
derived here. The proof relies on scalar white noise and a
conditioning pair of qubits. It makes no extension to arbitrary
mixed-unitary noise, higher dimensions, nontrivial Hamiltonians,
continuous-time dynamics or regularized communication capacities.
Neither the numerical gate checks nor the entropy sign should be
interpreted as an entanglement-breaking criterion.

\section{Methods and reproducibility}
The result is exact and finite-dimensional. Its error control is
the pair of analytical inequalities in the proof, with no
perturbative, finite-size or floating-point approximation.
The accompanying standard-library Python scripts reconstruct
the Choi input in two tensor conventions and test special-gate
certificates with rational arithmetic. A separate exact
complex-rational check covers eight unitary gates and four noise
weights, including rank-deficient and complex instances, with
32 primal, dual and marginal certificates and 128 sign-invariance
checks. The finite checks diagnose normalization and implementation
errors; they do not establish the universal quantifiers.
Regeneration instructions and the tested environment accompany the
code. The self-contained proof is sufficient for the analytic result.

Generative AI was used in proposing and developing the argument,
in writing reconstruction and certificate-checking code, in
literature searches and comparison, in automated scientific
checking, and in drafting this manuscript. These checks do not
constitute human verification or external peer review.

\paragraph{Acknowledgments.}
The research and drafting workflow used OpenAI Codex, including
its language-model assistance for the tasks described above.
Human author review and confirmation of the manuscript remain pending.

\paragraph{Data availability.}
No experimental data were used. The reproduction and exact
certificate scripts, fixed test inputs and regeneration instructions
are included in the accompanying source package.

\bibliographystyle{unsrt}
\bibliography{references}
\end{document}
